%% ============================================================
%% shared/services/compute/containers.tex
%% Containers on AWS — ECS, EKS, Fargate, ECR — SAP (some SAA)
%% ============================================================

\servicesection{Containers on AWS}{\saptag}{ECS, EKS, Fargate, ECR --- container orchestration}

\begin{keyfacts}
\begin{itemize}
  \item \textbf{Container}: packages application code + runtime + dependencies in a portable unit.
        Shares host OS kernel (lighter than VMs); starts in seconds.
  \item \textbf{Container image}: immutable snapshot; stored in a registry (ECR).
  \item \textbf{Orchestration}: manages scheduling, scaling, networking, and health of containers
        across a cluster.
\end{itemize}
\end{keyfacts}

\subsection{Container Platform Selection}

\begin{comparebox}[title={Container service comparison}]
\begin{tabular}{@{}p{3.2cm}p{5cm}p{3.8cm}@{}}
\toprule
\textbf{Service} & \textbf{Description} & \textbf{Choose when\ldots} \\
\midrule
\textbf{Amazon ECS}          & AWS-native container orchestration; simpler than Kubernetes;
                               tight AWS service integration &
  You want AWS-native simplicity; no Kubernetes expertise needed \\[4pt]
\textbf{Amazon EKS}          & Managed Kubernetes; AWS handles control plane &
  You need Kubernetes API compatibility, existing K8s tooling, or multi-cloud portability \\[4pt]
\textbf{AWS Fargate}         & Serverless compute engine for ECS or EKS; no EC2 instances to manage;
                               pay per vCPU/memory/second &
  You want to run containers without managing servers \\[4pt]
\textbf{Amazon ECR}          & Managed container image registry; integrates with ECS, EKS, Lambda,
                               CodeBuild; OCI-compliant &
  Always use ECR for images deployed to AWS services \\[4pt]
\textbf{ECS Anywhere}        & Run ECS tasks on on-premises servers using SSM agent &
  Hybrid; on-prem workloads with AWS control plane \\[4pt]
\textbf{EKS Anywhere}        & Run EKS on on-premises hardware (bare metal / VMware) &
  On-prem Kubernetes with EKS-compatible APIs \\[4pt]
\textbf{AWS App Runner}      & Fully managed; deploy container or source code; auto-scales;
                               no infrastructure knowledge needed &
  Simple web apps/APIs with no K8s/ECS expertise \\
\bottomrule
\end{tabular}
\end{comparebox}

\begin{examtip}
\textbf{ECS vs EKS:} ECS = simpler, AWS-native, less overhead. EKS = Kubernetes-compatible,
more complex, use when K8s is a hard requirement (team expertise, multi-cloud, K8s ecosystem tools).
\textbf{EC2 launch type vs Fargate:} EC2 = more control, GPU support, lower cost at sustained load.
Fargate = serverless, no cluster management, pay-per-task (better for variable/bursty workloads).
\end{examtip}

\subsection{Amazon ECS}

\begin{itemize}
  \item \textbf{Cluster}: logical grouping of tasks/services (EC2 or Fargate).
  \item \textbf{Task Definition}: blueprint specifying container image(s), CPU, memory, IAM role,
        networking mode, environment variables, log driver.
  \item \textbf{Task}: running instance of a task definition (one-time job).
  \item \textbf{Service}: maintains N running copies of a task; integrates with ALB/NLB;
        handles rolling deployments and health checks.
  \item \textbf{Task IAM Role}: credentials for the container (fine-grained, per-task; prefer over
        EC2 instance role for containers).
  \item \textbf{Networking modes}: \texttt{awsvpc} (each task gets own ENI + security group ---
        recommended), \texttt{bridge} (Docker NAT), \texttt{host}.
\end{itemize}

\begin{sapdepth}
\textbf{ECS blue/green with CodeDeploy:}
\begin{itemize}
  \item CodeDeploy registers new task set; shifts ALB traffic from blue (old) to green (new) target group.
  \item Original task set retained for configurable rollback window (e.g.\ 1 hour).
  \item Automatic rollback triggers on CloudWatch alarm.
  \item Supports canary (10\%$\rightarrow$100\%) or linear traffic shifting.
\end{itemize}
\end{sapdepth}

\subsection{Amazon EKS}

\begin{itemize}
  \item Managed Kubernetes control plane (API server, etcd); AWS patches and scales it.
  \item \textbf{Node groups}: managed EC2 instances (Auto Scaling Group); or Fargate profiles
        for serverless nodes.
  \item \textbf{EKS Add-ons}: managed VPC CNI, CoreDNS, kube-proxy, EBS CSI driver.
  \item \textbf{IRSA (IAM Roles for Service Accounts)}: map K8s service accounts to IAM roles;
        fine-grained pod-level AWS permissions.
  \item \textbf{Karpenter}: open-source node provisioner for EKS; replaces Cluster Autoscaler;
        faster node provisioning, better bin-packing.
\end{itemize}

\subsection{Amazon ECR}

\begin{itemize}
  \item Fully managed Docker-compatible container image registry.
  \item \textbf{Private repositories}: access via IAM (resource-based policies for cross-account).
  \item \textbf{Public Gallery}: host public images (\texttt{public.ecr.aws}).
  \item \textbf{Image scanning}: on-push or continuous scanning via Inspector; reports CVEs.
  \item \textbf{Lifecycle policies}: automatically delete untagged or old images to save storage cost.
  \item \textbf{Replication}: replicate images across regions and accounts for multi-region deployments.
\end{itemize}

\subsection{Compute Platform Selection (Full Spectrum)}

\begin{comparebox}[title={Compute platform decision guide}]
\begin{tabular}{@{}p{3.5cm}p{8cm}@{}}
\toprule
\textbf{Platform} & \textbf{Choose when\ldots} \\
\midrule
\textbf{EC2}              & Full control needed; persistent workloads; custom OS; GPU; lift-and-shift \\
\textbf{Elastic Beanstalk} & PaaS; auto-manages infra (EC2, ALB, ASG); web apps and APIs; replatform \\
\textbf{Lambda}           & Event-driven; short-duration; serverless; auto-scales to zero \\
\textbf{ECS on Fargate}   & Containers; no cluster management; variable load \\
\textbf{ECS on EC2}       & Containers; sustained high load; GPU workloads; cost optimisation \\
\textbf{EKS}              & Kubernetes required; multi-cloud; K8s ecosystem tooling \\
\textbf{App Runner}       & Simple web apps; source code or container; no infra knowledge \\
\textbf{AWS Outposts}     & On-premises; data residency; hybrid; low-latency to on-prem systems \\
\bottomrule
\end{tabular}
\end{comparebox}
